 

\section{Theoretical framework}
\label{2405:sec:app:framework}

\subsection{Difference inequalities for vector-valued processes}

We begin with a lemma that generalizes \cite[Proposition 4.1]{arous2021online}. Let $\Phi: \R^d \to \R^d$ be a random function, and consider the process $(\vec{u}_t)_{t \geq 0} \in \mathbb{S}^{d-1}$ given by
\[ \vec{u}_{t+1} = \frac{\vec{u}_t + \gamma \Phi(\vec{u}_t)}{\norm{\vec{u}_t + \gamma \Phi(\vec{u}_t)}}. \]
We make the following assumptions:
\begin{assumption}\label{2405:assump:app:initialization}
  The initial value of $\vec{u_0} \in \mathbb{S}^{d-1}$ is approximately isotropic, in the sense that there exists a constant $c > 0$ such that
  \[ \Ea{\vec{u_0}\vec{u_0}^\top} \succ \frac{c}d I\]
\end{assumption}
\begin{assumption}\label{2405:assump:app:concentration}
  There exist constants $C, \iota > 0$ such that
  \begin{align}
    \sup_{\vec{u} \in \mathbb{S}^{d-1}}  \norm{\Ea{\Phi(\vec u)\Phi(\vec u)^\top}} &\leq C \\
    \sup_{\vec{u} \in \mathbb{S}^{d-1}}  \Ea{\norm{\Phi(\vec u)}^{4 + \iota}} &\leq C N^{\frac{4 + \iota}2}
  \end{align}
  Further, we have $\langle \vec{u}, \Phi(\vec{u}) \rangle = 0$ a.s. for any $\vec{u} \in \mathbb{S}^{d-1}$.
\end{assumption}

We fix a vector $\vec{v}$ and consider the random process $m_t(\vec{v}) := \langle \vec{u}_t, \vec{v} \rangle$. For any $s \geq 0$, we define the stopping times
\[ \tau_\theta^+(\vec{v}) = \inf \{t \geq 0 : m_t(\vec{v}) \geq \theta \} \quad \text{and} \quad \tau_\theta^-(\vec{v}) = \inf \{t \geq 0 : m_t(\vec{v}) \leq \theta \}  \]
\begin{proposition}\label{2405:prop:app:diff_equation}
  Fix an $\alpha > 0$. There exist constants $c(\delta, \alpha), \kappa(\delta) > 0$ such that if
  \begin{equation}\label{2405:eq:app:bound_gamma_T}
    \gamma \leq c(\delta, \alpha)d^{-1} \quad \text{and} \quad T \leq c(\delta, \alpha) \gamma^{-2} d^{-1},
  \end{equation}
  then, with probability $1 - \delta$, the following holds. The initialization $m_0(\vec{v})$ satisfies
  \begin{equation}
    |m_0(\vec{v})| \geq \frac{\kappa}{\sqrt{d}},
  \end{equation}
  and:
  \begin{itemize}
    \item if $m_0(v) > 0$, for any $t \leq T \wedge \tau_0^-(\vec{v})$,
      \begin{equation}\label{2405:eq:app:diff_equation_positive}
        (1 - \alpha) m_0(\vec{v}) + \gamma\sum_{s=0}^{t-1} \langle \Psi^+(\vec{u}_s), \vec{v} \rangle \leq m_t(\vec{v}) \leq (1 + \alpha) m_0(\vec{v}) + \gamma\sum_{s=0}^{t-1} \langle \Psi(\vec{u}_s), \vec{v} \rangle,
      \end{equation}
    \item if $m_0(v) < 0$, for any $t \leq T \wedge \tau_0^+(\vec{v})$,
      \begin{equation}\label{2405:eq:app:diff_equation_negative}
        (1 + \alpha) m_0(\vec{v}) + \gamma\sum_{s=0}^{t-1} \langle \Psi(\vec{u}_s), \vec{v} \rangle \leq m_t(\vec{v}) \leq (1 - \alpha) m_0(\vec{v}) + \gamma\sum_{s=0}^{t-1} \langle \Psi^-(\vec{u}_s), \vec{v} \rangle,
      \end{equation}
  \end{itemize}
  where $\Psi, \Psi^{\pm}$ is a deterministic function equal to
  \[ \Psi(\vec{u}) = \Ea{\Phi(\vec{u})} \quad \text{and} \quad \Psi^{\pm}(\vec{u}) = \Psi(\vec{u}) \mp \gamma \Ea{\norm{\Phi(\vec{u})}^2}\vec{u}.\]
\end{proposition}

\begin{proof}
  We assume for simplicity that $m_0(\vec{v}) > 0$; the case where $m_0(\vec{v}) > 0$ is treated identically. From Assumption~\ref{2405:assump:app:initialization}, by a simple Markov inequality, there exists a constant $\kappa(\delta) > 0$ such that
  \begin{equation} \label{2405:eq:app:large_init_event}
    \dP\left(|m_0(\vec{v})| \geq \frac{\kappa}{\sqrt{d}}\right) \geq 1 - \frac{\delta}3.
  \end{equation}
  The update equation for $m_t(\vec{v})$  reads
  \[ m_{t+1}(\vec{v}) = \frac{m_t(\vec{v}) + \gamma \langle \Phi(\vec{u}_t), \vec{v} \rangle}{\sqrt{1 + \gamma^2  \norm{\Phi(\vec{u}_t)}^2}}. \]
  We then have two cases:
  \begin{itemize}
    \item either $m_t(\vec{v}) + \gamma \langle \Phi(\vec{u}_t), \vec{v} \rangle < 0$, which implies that $t+1 = \tau_0^-(\vec{v})$ and there is nothing to prove,
    \item or we can use the inequality $\frac1{\sqrt{1+x}} \geq 1 - x$, which implies that
      \begin{align*}
        m_{t+1}(\vec{v}) &\geq (m_t(\vec{v}) + \gamma \langle \Phi(\vec{u}_t), \vec{v} \rangle)(1 - \gamma^2  \norm{\Phi(u_t)}^2) \\
        &= (1 - \gamma^2  \norm{\Phi(u_t)}^2)m_t + \gamma \langle \Phi(\vec{u}_t), \vec{v} \rangle - \gamma^3 \langle \Phi(\vec{u}_t), \vec{v} \rangle  \norm{\Phi(u_t)}^2
      \end{align*}
  \end{itemize}
  By recursion, deterministically,
  \[
    \resizebox{\linewidth}{!}{$\displaystyle
      m_t(\vec{v}) \geq m_0(\vec{v}) + \gamma\sum_{s = 0}^{t-1} \left(\langle \Phi(\vec{u}_s), \vec{v} \rangle - \gamma\norm{\Phi(\vec{u}_s)}^2\langle \vec{u}_s, \vec{v} \rangle\right) - \sum_{s=0}^{t-1}\gamma^3 \left|\langle \Phi(\vec{u}_s), \vec{v} \rangle\right|\norm{\Phi(\vec{u}_s)}^2
    $}
  \]

  Since Assumption~\ref{2405:assump:app:concentration} implies Assumption B in \cite{arous2021online}, for a small enough choice of $c(\delta, \alpha)$, one has
  \[ \dP\left( \sum_{t=0}^T\gamma^3 \left|\langle \Phi(\vec{u}_t), \vec{v} \rangle\right| \norm{\Phi(\vec{u}_s)}^2 \leq \frac{\alpha\kappa}{2\sqrt{d}} \right) \geq 1 - \frac\delta6. \]

  Additionnally, the same assumption implies that
  \[
    \resizebox{\linewidth}{!}{$\displaystyle
      \operatorname{Var}(\langle \Phi(\vec{u}_s), \vec{v} \rangle)  \leq \Ea{\langle \Phi(\vec{u}_s), \vec{v} \rangle^2} \leq C_1 \quad \text{and} \quad \operatorname{Var}(\gamma\norm{\Phi(\vec{u}_s)}^2) \leq \gamma^2\Ea{\norm{\Phi(\vec{u}_s)}^4} \leq c(\delta)^2 C_1.
    $}
  \]
  We can then apply Doob's martingale inequality: there exists a constant $C(\delta)$ such that with probability at least $1 - \frac\delta6$
  \[
    \resizebox{\linewidth}{!}{$\displaystyle
      \sup_{t \leq T} \left| \gamma\sum_{s = 0}^t  \left( \langle \Phi(\vec{u}_s), \vec{v} \rangle - \gamma\norm{\Phi(\vec{u}_s)}^2\langle \vec{u}_s, \vec{v} \rangle - \Ea{\langle \Phi(\vec{u}_s), \vec{v} \rangle - \gamma\norm{\Phi(\vec{u}_s)}^2\langle \vec{u}_s, \vec{v} \rangle} \right) \right| \leq C(\delta)\gamma\sqrt{T}.
    $}
  \]
  The bound~\eqref{2405:eq:app:bound_gamma_T} on $T$ implies that we can choose $c(\delta, \alpha)$ small enough so that
  \[ C(\delta)\gamma\sqrt{T} \leq \frac{C(\delta) \sqrt{c(\delta, \alpha)}}{\sqrt{d}} \leq \frac{\alpha\kappa}{2\sqrt{d}}.\]
  Combining all those inequalities with a union bound, under the event of Eq.~\eqref{2405:eq:app:large_init_event},
  \begin{equation}\label{2405:eq:app:diff_equation_lower}
    \dP\left( m_t(\vec{v}) \geq (1 - \alpha) m_0(\vec{v}) + \gamma\sum_{s=0}^{t-1} \langle \Psi^+(\vec{u}_s), \vec{v} \rangle \quad \forall t < T \wedge \tau_0^-(\vec{v})\right) \geq 1 - \frac\delta3,
  \end{equation}
  having used that $\langle \Psi^+(\vec{u}), \vec{v} \rangle = \Ea{\langle \Phi(\vec{u}), \vec{v} \rangle - \gamma\norm{\Phi(\vec{u})}^2\langle \vec{u}, \vec{v} \rangle}$.

  The upper bound follows similarly:  using this time the fact that for $x > 0$ we have $\frac1{\sqrt{1+x}} \leq 1$, we have under \eqref{2405:eq:app:large_init_event}
  \begin{equation}\label{2405:eq:app:diff_equation_upper}
    \dP\left( m_t(\vec{v}) \leq (1 + \alpha) m_0(\vec{v}) + \gamma\sum_{s=0}^{t-1} \langle \Psi(\vec{u}_s), \vec{v} \rangle \quad \forall t < T \wedge \tau_0^-(\vec{v})\right) \geq 1 - \frac\delta3.
  \end{equation}
  A union bound between the events of Eqs.~\eqref{2405:eq:app:large_init_event}, \eqref{2405:eq:app:diff_equation_lower}, \eqref{2405:eq:app:diff_equation_upper} concludes the proof.
\end{proof}

\subsection{Low-dimensional sufficient statistics}

In \cite{arous2021online}, the function $\Psi(\vec{u})$ can be written as $\Psi(\vec{u}) = \phi(\langle \vec{u}, \vec{v} \rangle)\vec{v}$. As a result, the bounds of Proposition~\ref{2405:prop:app:diff_equation} are actually a closed function of $m_t(\vec{v})$. We generalize this phenomenon to a larger class:

\begin{assumption}\label{2405:assump:app:low_dim}
  There exists a $q > 0$ such that the function $\Psi$ is $q$-approximately $k$-dimensional, in the sense that there exist a function $\psi: \R^k \to \R^k$ and an orthonormal matrix $W^\star \in \R^{k \times d}$ such that for any $\vec{u}\in\R^d$,
  \[ \norm{ W^\star \Psi(\vec{u}) - \psi(W^\star \vec{u})} \leq \left( \norm{W^\star \vec{u}} \vee \frac1{\sqrt{d}} \right) \cdot q. \]
\end{assumption}

By an abuse of notation, we shall sometimes view $\psi$ as a function from $V^\star$ to itself, where $V^\star$ is the span of the columns of $W^\star$. However, Assumption~\ref{2405:assump:app:low_dim} ensures that all quantities related to $\psi$ are only depending on the (oftentimes fixed) dimension $k$ instead of $d$.

It turns out that in this case, we can control accurately the trajectory of $W^\star \vec{u}$, at least when $f$ is not too flat around $\vec{0}$. Let $\Pi^\star = (W^{\star})^{\top}W^\star$ be the projection operator on $V^\star$, and define the stopping time
\[ \tau_\theta = \min \{t \geq 0 : \norm{\Pi^\star \vec{u}} \geq \theta \}. \]

\begin{proposition}\label{2405:prop:app:low_dim_linear}
  Assume that Condition \ref{2405:assump:app:low_dim} holds with $q$ small enough, and that $\langle \vec{u}_0, \psi(\vec{0}) \rangle > 0$ (in particular, $\psi(\vec{0}) \neq 0$). For any $\delta, \epsilon > 0$, there exist constants $c(\delta, \epsilon), \eta(\epsilon)$ such that if
  \begin{equation}
    \gamma \leq c d^{-1} \quad \text{and} \quad T \leq c  \gamma^{-2} d^{-1},
  \end{equation}
  then
  \begin{equation}
    \dP\left(\norm{\Pi^\star \vec{u}_t - \gamma t \psi(\vec{0})} \leq \sqrt{k} \eta \epsilon \quad \forall t \leq T \wedge \tau_\eta \right) \geq 1 - \delta.
  \end{equation}
\end{proposition}

\begin{proof}
  Let $\vec{v}_1, \dots, \vec{v}_k$ be an orthonormal basis of $V^\star$, with $\vec{v}_1$ aligned with $\psi(\vec{0})$. By applying Proposition~\ref{2405:prop:app:diff_equation} to $\delta' = \frac{\delta}{2k}$ and $\alpha = \frac12$, with probability $1 - \frac\delta2$, for any $i \in [k]$ and $t \leq T \wedge \tau^{\pm}(\vec{v}_i)$, $m_t(\vec{v}_i)$ satisfies the inequality~\eqref{2405:eq:app:diff_equation_positive} or \eqref{2405:eq:app:diff_equation_negative} (depending on the sign of $m_0(\vec{v_i}))$. In particular, for $\vec{v}_1 = \psi(\vec{0})$, since we assumed that the initial sign is positive, we have
  \begin{equation}\label{2405:eq:app:u_psi_ineq}
    \frac12 \langle \vec{u}_0, \psi(\vec{0}) \rangle + \gamma\sum_{s=0}^{t-1} \langle \Psi^+(\vec{u}_s), \psi(\vec{0}) \rangle \leq \langle \vec{u}_t, \psi(\vec{0}) \rangle \leq  \frac32 \langle \vec{u}_0, \psi(\vec{0}) \rangle + \gamma\sum_{s=0}^{t-1} \langle \Psi(\vec{u}_s), \psi(\vec{0}) \rangle.
  \end{equation}

  Define $C(\delta)$ such that  $\langle \vec{u}_t, \psi(\vec{0}) \rangle \leq C(\delta)d^{-1/2}$ with probability at least $1 - \frac{\delta}{2k}$. By a Taylor expansion, there exists an $\eta(\epsilon) > 0$ such that for any $\vec{v} \in \R^d$ with $\norm{\vec{v}} \leq \eta$, $\norm{\psi(\vec{v}) - \psi(\vec{0})} \leq \epsilon'$, where $\epsilon'$ will be fixed later. With Assumption~\ref{2405:assump:app:low_dim}, this implies that for any $\vec{u} \in \R^d$,
  \[ \norm{\psi(\vec{0})}^2 - (\epsilon' + \eta q) \norm{\psi(\vec{0})} \leq \langle \Psi(\vec{u}), \psi(\vec{0}) \rangle \leq \norm{\psi(\vec{0})}^2 + (\epsilon' + \eta q) \norm{\psi(\vec{0})}. \]
  Further, by Jensen's inequality, $\Ea{\norm{\Phi(\vec{u})}^2} \leq C_1 d$ so by choosing $c(\delta, \epsilon)$ small enough, we can ensure that
  \[ \langle \Psi^+(\vec{u}), \psi(\vec{0}) \rangle \geq \norm{\psi(\vec{0})}^2 - (2\epsilon + \eta q) \norm{\psi(\vec{0})} \]
  Plugging those estimates in \eqref{2405:eq:app:u_psi_ineq}, we find that for $t \leq T \wedge \tau_0^-(\vec{v_1}) \wedge \tau_\eta$, and $d$ large enough so that
  \[ C(\delta)d^{-1/2}\leq \frac{\eta\epsilon}{3} \norm{\psi(\vec{0})}, \]
  we have
  \begin{equation}
    \resizebox{\linewidth}{!}{$\displaystyle
      \gamma t \left(\norm{\psi(\vec{0})}^2 - (2\epsilon' + \eta q) \norm{\psi(\vec{0})}\right)  \leq \langle \vec{u}_t, \psi(\vec{0}) \rangle \leq \frac{\eta\epsilon}{2} \norm{\psi(\vec{0})} + \gamma t  \left(\norm{\psi(\vec{0})}^2 + (\epsilon' + \eta q) \norm{\psi(\vec{0})}\right).
    $}
  \end{equation}

  For small enough $\epsilon$ and $q$, the LHS of the above inequality is always at least $\gamma t \norm{\psi(0)}^2 / 2$, which implies that
  \begin{equation}\label{2405:eq:app:bound_tau_eta}
    \tau_0^-(\vec{v}_1) \geq \tau_\eta \quad \text{and} \quad T \wedge \tau_\eta \leq \frac{2 \eta}{\gamma \norm{\psi(\vec{0})}}
  \end{equation}
  Choosing $\epsilon'$ and $\eta$ small enough so that
  \[ \frac{(2\epsilon' + \eta q)}{\norm{\psi(\vec{0})}^2} \leq \frac{\epsilon}{2}\]
  we ensure that for $t \leq T \wedge \tau_\eta$ we have
  \begin{equation}
    \left| \langle \vec{u_t} - \gamma t \psi(\vec{0}), \psi(\vec{0}) \rangle \right| \leq \epsilon \eta \norm{\psi(\vec{0})}.
  \end{equation}

  We now move on to the other $\vec{v}_i$. For $i \geq 2$, $\vec{v}_i$ is orthogonal to $\psi(\vec{0})$, and this time Assumption~\ref{2405:assump:app:low_dim} implies
  \[ \left| \langle \Psi(\vec{u}), \vec{v}_i \rangle \right| \leq \epsilon + \eta q. \]
  Regardless of the sign of $m_0(\vec{v}_i)$, this implies that for $t \leq T \wedge \tau_\eta$,
  \[ |\langle \vec{u_t}, \vec{v}_i \rangle| \leq \frac32|\langle \vec{u}_0, \vec{v}_i \rangle| + \gamma t (\epsilon' + \eta q) \leq \frac{\eta\epsilon}{2} + \frac{2\eta(2\epsilon' + \eta q)}{\norm{\psi(\vec{0})}}, \]
  having used the bound~\eqref{2405:eq:app:bound_tau_eta} on $\tau_\eta$ and bounding the initial value as above with probability $1- \frac\delta{2k}$. We can always choose $\eta, \epsilon'$ such that the above bound is at most $\epsilon$; since
  \[ \norm{\Pi_V \vec{u}_t - \gamma t \psi(\vec{0})}^2 = \frac{\langle \vec{u_t} - \gamma t \psi(\vec{0}), \psi(\vec{0}) \rangle^2}{\norm{\psi(\vec{0})}^2} + \sum_{i=2}^k \langle \vec{u_t}, \vec{v}_i \rangle^2 \leq k\epsilon^2, \]
  this concludes the proof.
\end{proof}

\begin{proposition}\label{2405:prop:app:low_dim_exp}
  Assume that $\psi(\vec{0}) = \vec{0}$, and $H_\psi := d\psi(\vec{0}) \neq 0$ is symmetric. Let $\bar\lambda$ be the largest eigenvalue of $H_\psi$, and assume that $\bar\lambda > 0$.
  For any $\delta, \epsilon > 0$, there exist constants $c(\delta, \epsilon), \eta(\epsilon), c'(\delta, \epsilon), C(\delta, \epsilon)$ such that if
  \begin{equation}
    \gamma \leq c d^{-1} \quad \text{and} \quad T \leq c  \gamma^{-2} d^{-1},
  \end{equation}
  then
  \begin{equation}
    \dP\left(\norm{\Pi^\star \vec{u}_t} \geq c' \norm{\Pi^\star \vec{u}_0} e^{\gamma(\bar\lambda - Ck(\epsilon + q))t} \quad \forall t \leq T \wedge \tau_\eta \right) \geq 1 - \delta.
  \end{equation}
  Further, the mass of $\Pi^\star \vec{u}_t$ is concentrated on the top eigenspace of $H_\psi$: there exists a $\tilde\lambda < \bar\lambda - Ck(\epsilon + q)$ such that on the same event,
  \[ \norm{\Pi^\star \vec{u}_{\tau_\eta} - \bar\Pi^\star \vec{u}_t} \leq \norm{\Pi^\star \vec{u}_0} e^{\gamma \tilde\lambda t} \quad \forall t \leq T \wedge t_\eta, \]
  where $\bar\Pi^\star$ is the projection on the associated to $\bar\lambda$.
\end{proposition}

\begin{proof}
  Fix $(\vec{v}_1, \dots, \vec{v}_k)$ an orthonormal basis of $V^\star$, consisting of eigenvectors of $H_\psi$, and let $\lambda_i$ be the eigenvalue associated with $\vec{v}_i$. The largest eigenvalue of $H_\psi$ that differs from $\bar\lambda$ will be denoted $\underline\lambda$ (we can let $\underline\lambda = -\infty$ if $H_\psi = \bar\lambda I$).
  
  We choose an $\eta(\epsilon)$ such that for $\vec{v}\in \R^k$ such that $\norm{\vec{v}} \leq \eta$,
  \[ \norm{\psi(\vec{v}) - H_\psi \vec{v}} \leq \epsilon \norm{\vec{v}}. \]
  In particular, for any $i \in [k]$ and any vector $\vec{u}$, we have for small enough choice of $c(\delta, \epsilon)$
  \[ \langle \Psi^+(\vec{u}), \vec{v}_i \rangle \geq \lambda_i \langle \vec{v}_i, \vec{u} \rangle - (2\epsilon + q)\norm{\Pi^\star\vec{u}} \quad \langle \Psi(\vec{u}), \vec{v}_i \rangle \leq \lambda_i \langle \vec{v}_i, \vec{u} \rangle + (\epsilon + q)\norm{\Pi^\star\vec{u}} \]
  Applying Proposition~\ref{2405:prop:app:diff_equation} to $\alpha = \epsilon$, and assuming that $m_0(\vec{v}_i) > 0$, we find that on an event with probability $1 - \delta/(2k)$, for $t \leq T \wedge \tau_0^+(\vec{v_i}) \wedge \tau_\eta$,
  \begin{align*}
     \langle \vec{v}_i, \vec{u}_t \rangle &\geq (1 - \epsilon)m_0(\vec{v}_i) + \gamma \sum_{s=0}^{t-1} \left(\lambda_i\langle \vec{v}_i, \vec{u}_s \rangle  - (2\epsilon + q)\norm{\Pi^\star\vec{u}_s}\right) \\ 
     \langle \vec{v}_i, \vec{u}_t \rangle &\leq (1 + \epsilon)m_0(\vec{v}_i) + \gamma \sum_{s=0}^{t-1} \left(\lambda_i\langle \vec{v}_i, \vec{u}_s \rangle + (\epsilon + q)\norm{\Pi^\star\vec{u}_s}\right)
  \end{align*}
  The same argument when $m_0(\vec{v}_i) < 0$, we obtain
  that with probability $1 - \delta/2$, for any $t \leq T \wedge \tau_\eta$ and $i \in [k]$:
  \begin{align}
    |\langle \vec{v}_i, \vec{u}_t \rangle| &\geq (1 - \epsilon)|\langle \vec{v}_i, \vec{u}_0 \rangle| + \gamma \sum_{s=0}^{t-1} \left(\lambda_i|\langle \vec{v}_i, \vec{u}_s \rangle|  - (2\epsilon + q)\norm{\Pi^\star\vec{u}_s}\right) \label{2405:eq:app:lower_bound_exp}\\
    |\langle \vec{v}_i, \vec{u}_t \rangle| &\leq (1 + \epsilon)|\langle \vec{v}_i, \vec{u}_0 \rangle| + \gamma \sum_{s=0}^{t-1} \left(\lambda_i|\langle \vec{v}_i, \vec{u}_s \rangle| + (\epsilon + q)\norm{\Pi^\star\vec{u}_s}\right)\label{2405:eq:app:upper_bound_exp}
  \end{align}
Define $\bar\Pi^\star$ the projection on the eigenspace associated to $\bar\lambda$, and $\underline\Pi^\star$ its orthogonal projection. Summing the upper bound of \eqref{2405:eq:app:upper_bound_exp} on all $v_i$ with $\lambda_i \leq \underline\lambda$,
\[ \norm{\underline\Pi^\star \vec{u}_t}_1 \leq (1 + \epsilon)\norm{\underline\Pi^\star \vec{u}_0}_1 + \gamma (\underline\lambda + k(\epsilon + q)) \sum_{s=0}^{t-1} \norm{\Pi^\star \vec{u}_s}_1, \]
  where the $\ell^1$ norm $\norm{ \cdot }_1$ in $V^\star$ is computed w.r.t the basis $(\vec{v}_1, \dots, \vec{v}_k)$, and we used the norm bound $\norm{\Pi^\star \vec{u}_s} \leq \norm{\Pi^\star \vec{u}_s}_1$.

  Call $\cP(t_0)$ the following condition:
  \begin{equation}\label{2405:eq:app:rec_condition}
      \norm{\underline\Pi^\star \vec{u}_t}_1 \leq C\norm{\bar\Pi^\star \vec{u}_t}_1 \quad \text{for all} \  t < t_0, ,
  \end{equation} where $C$ will be fixed later, and assume that $\cP(t_0)$ holds for some $t_0 > 0$.
  Using the discrete Grönwall inequality, the following holds: for any $t \leq T \wedge \tau_\eta \wedge t_0$,
  \begin{equation}\label{2405:eq:app:bad_upper_bound}
      \norm{\underline\Pi^\star \vec{u}_t}_1 \leq (1 + \epsilon)\norm{\underline\Pi^\star \vec{u}_0}_1 e^{\gamma (\underline\lambda + (1+C)k(\epsilon + q))t}.
  \end{equation} 
   We now sum the lower bound of \eqref{2405:eq:app:lower_bound_exp} on all $v_i$ with $\lambda_i = \bar\lambda$, to obtain
    \[ \norm{\bar\Pi^\star\vec{u}_t}_1 \geq (1 - \epsilon) \norm{\bar\Pi^\star \vec{u}_0}_1 + \gamma \sum_{s=0}^{t-1} (\bar\lambda - k(2\epsilon + q)) \norm{\bar\Pi^\star \vec{u}_s}_1 - k(2\epsilon + q)\norm{\underline\Pi^\star \vec{u}_s}_1. \]
      above bound implies that for all $t \leq t_0$,
     \[ \norm{\bar\Pi^\star\vec{u}_t}_1 \geq (1 - \epsilon) \norm{\bar\Pi^\star \vec{u}_0}_1 + \gamma \sum_{s=0}^{t-1} (\bar\lambda - (1+C)k(2\epsilon + q)) \norm{\bar\Pi^\star \vec{u}_s}_1. \]
     By the discrete Grönwall inequality,
     \begin{align} 
        \norm{\bar\Pi^\star\vec{u}_t}_1 &\geq (1 - \epsilon) \norm{\bar\Pi^\star \vec{u}_0}_1 \left(1 +  \gamma(\bar\lambda - (1+C)k(2\epsilon + q)) \right)^t \nonumber\\ 
        &\geq (1 - \epsilon) \norm{\bar\Pi^\star \vec{u}_0}_1 \exp\left(\gamma(\bar\lambda - (1+C)k(2\epsilon + q))t - O(\gamma^2 T) \right) \nonumber\\
        &\geq (1 - 2\epsilon) \norm{\bar\Pi^\star \vec{u}_0}_1 \exp\left(\gamma(\bar\lambda - (1+C)k(2\epsilon + q))t\right), \label{2405:eq:app:good_lower_bound}
     \end{align}
     where the last inequality holds for large enough $d$ because of the bound \eqref{2405:eq:app:bound_gamma_T}.
     Now, fix a constant $r(\delta)$ such that with probability $1 - \delta/2$, $\norm{\bar\Pi^\star \vec{u}_0}_1 \leq r(\delta)\norm{\underline\Pi^\star \vec{u}_0}_1$. Then condition $\cP(1)$ holds as long as $C > r(\delta)$, and $\cP(t_0)$ implies $\cP(t_0+1)$ whenever the lower bound of eq. \eqref{2405:eq:app:good_lower_bound} is higher than the upper bound of eq. \eqref{2405:eq:app:bad_upper_bound} for $t \leq t_0$. In particular, this holds whenever
     \[ C \geq \frac{1 + \epsilon}{1 - 2\epsilon}r(\delta) \quad \text{and} \quad \bar\lambda - (1+C)k(2\epsilon + q) \geq  \underline\lambda + (1+C)k(\epsilon + q).\]
     The above condition can be satisfied if and only if
     \[ \frac{1 + \epsilon}{1 - 2\epsilon}r(\delta) < \frac{\bar\lambda - \underline\lambda}{3\epsilon + 2q} - 1, \]
     which can always be satisfied when $\epsilon, q$ are low enough. The proposition ensues from the remark that $\norm{\bar\Pi^\star\vec{u}_t} \leq \norm{\Pi^\star \vec{u}} \leq (1+C)\norm{\bar\Pi^\star \vec{u}_t}$ for any $t \leq T \wedge \tau_\eta$.
\end{proof}

\subsection{Bounding the hitting times}

It finally remains to show how the bounds of Propositions \ref{2405:prop:app:low_dim_linear} and \ref{2405:prop:app:low_dim_exp} allow us to bound the hitting times $\tau_\eta$. We summarize those in the following proposition:
\begin{proposition}\label{2405:prop:app:hitting_times}
    Let $\vec{u}$ be a process that satisfies Assumptions \ref{2405:assump:app:initialization}-\ref{2405:assump:app:low_dim} for low enough $q$. Then for any $\delta, \epsilon > 0$, there exist constants $\gamma_0(\delta, \epsilon), \eta(\epsilon), \kappa(\delta), K(\delta), C(\delta, \epsilon)$ such that, with probability $1 - \delta$:
    \begin{enumerate}
        \item if $\psi(\vec{0}) \neq \vec{0}$ and $\langle \vec{u}_0, \psi(\vec{0}) \rangle > 0$, letting $\gamma = \gamma_0 d^{-1}$,
        \[ \frac{\eta(1 - \sqrt{k}\epsilon)}{\gamma \norm{\psi(\vec{0})}} \leq \tau_\eta \leq \frac{\eta(1 + \sqrt{k}\epsilon)}{\gamma \norm{\psi(\vec{0})}} \]
        \item if $\psi(\vec{0}) = 0$ and $d\psi(\vec{0})$ is symmetric and has a positive maximum eigenvalue $\bar\lambda$, then letting $\gamma = \gamma_0 (d\log(d))^{-1}$
        \[ \frac{\log\left(\frac{\eta\sqrt{d}}{(1+\epsilon)K}\right)}{\gamma(\bar\lambda + k(\epsilon + q))} \leq \tau_\eta \leq \frac{\log\left(\frac{\eta\sqrt{d}}{\kappa c'}\right)}{\gamma(\bar\lambda - Ck(\epsilon + q))}\]
    \end{enumerate}
    Further, the directions of $\Pi^\star \vec{u}_{\tau_\eta}$ concentrate, in the following sense:
    \begin{enumerate}
        \item if $\psi(\vec{0}) \neq \vec{0}$,
        \[ \norm{\frac{\Pi^\star \vec{u}_{\tau_\eta}}{\eta} -  \frac{\psi(\vec{0})}{\norm{\psi(\vec{0})}}} \leq 2\sqrt{k}\epsilon \]
        \item if $\psi(\vec{0}) = \vec{0}$, for some $\beta < 1$,
        \[ \norm{\Pi^\star \vec{u}_{\tau_\eta} - \bar \Pi^\star \vec{u}_{\tau_\eta}} \leq C(\delta, \epsilon) d^{-\frac{1 - \beta}{2}},\]
        where $\bar\Pi^\star$ is the projection on the top eigenspace of $d\psi(\vec{0})$.
    \end{enumerate}
\end{proposition}

\begin{proof}
    The proof for the linear case is almost immediate. From Proposition \ref{2405:prop:app:low_dim_linear}, with probability $1 - \delta$, we have for all $t \leq T \wedge \tau_\eta$
    \[ \gamma t \norm{\psi(\vec{0})} - \sqrt{k}\eta\epsilon \leq \norm{\Pi^\star\vec{u}_t} \leq \gamma t \norm{\psi(\vec{0})} + \sqrt{k}\eta\epsilon \]
    Define
    \[ \bar t_\epsilon = \frac{\eta(1 + \sqrt{k}\epsilon)}{\gamma \norm{\psi(\vec{0})}} \quad \text{and} \quad \underline t_\epsilon = \frac{\eta(1 - \sqrt{k}\epsilon)}{\gamma \norm{\psi(\vec{0})}}. \]
    For $t \leq \underline t_\epsilon$, the upper bound on $\norm{\Pi^\star \vec{u}_t}$ is lower than $\eta$, which directly means that $\tau_\eta \geq \underline t_\epsilon$. The lower bound is higher than $\eta$ for $t = \bar t_\epsilon$, which implies that $\tau_\eta \leq \bar t_\epsilon$ as long as $T \geq \bar t_\epsilon$. Such a $T$ satisfies the bound \eqref{2405:eq:app:bound_gamma_T} as long as
    \[ \frac{\eta(1 + \sqrt{k}\epsilon)}{\gamma \norm{\psi(\vec{0})}} \leq c(\delta, \epsilon) \gamma^{-2}d^{-1}, \]
    which is equivalent to
    \[ \gamma \leq \frac{c(\delta,\epsilon)\norm{\psi(\vec{0})}}{\eta(1 + \sqrt{k}\epsilon)}  d^{-1},\]
    which can always be ensured by decreasing $\gamma_0(\delta, \epsilon)$. Finally, we use the bound
    \[ \norm{\frac{\vec{x}}{\norm{\vec{x}}} - \frac{\vec{y}}{\norm{\vec{y}}}} \leq 2 \frac{\norm{\vec{x} - \vec{y}}}{\norm{\vec{x}}} \]
    to conclude that
    \[ \norm{\frac{\Pi^\star \vec{u}_{\tau_\eta}}{\eta} - \frac{\psi(\vec{0})}{\norm{\psi(\vec{0})}}} \leq 2\sqrt{k}\epsilon. \]

    We now move on to the quadratic case. By Proposition \ref{2405:prop:app:low_dim_exp} and the same reasoning as in the proof of \eqref{2405:eq:app:bad_upper_bound}, we have with probability $1 - \delta$, for all $t \leq T \wedge \tau_\eta$,
    \[ c' \norm{\Pi^\star \vec{u}_0} e^{\gamma(\bar\lambda - Ck(\epsilon + q))t} \leq \norm{\Pi^\star \vec{u}_t} \leq (1 + \epsilon)\norm{\Pi^\star \vec{u}_0} e^{\gamma(\bar\lambda + k(\epsilon + q))t} \]
    Upon redefining $\delta$, we can assume that the ``good'' event implies that
    \[ \frac{\kappa(\delta)}{\sqrt{d}} \leq \norm{\Pi^\star \vec{u}_0} \leq \frac{K(\delta)}{\sqrt{d}} \]
    for constants $k, K$. As in the linear case, the upper bound above implies that
    \[ \tau_\eta \geq \frac{\log\left(\frac{\eta\sqrt{d}}{(1+\epsilon)K}\right)}{\gamma(\bar\lambda + k(\epsilon + q))}, \]
    while the lower bound implies that
    \[ T \wedge \tau_\eta \leq \frac{\log\left(\frac{\eta\sqrt{d}}{\kappa c'}\right)}{\gamma(\bar\lambda - Ck(\epsilon + q))} =: \bar t_\epsilon'. \]
    It therefore remains to show that $\bar t_\epsilon' \leq T$; for large enough $d$ and small enough $\epsilon, q$, we can write $\bar t_\epsilon' \leq C'(\delta, \epsilon) \gamma^{-1} \log(d)$, and
    \[ C'(\delta, \epsilon) \gamma^{-1} \log(d) \leq c\gamma^{-2}d^{-1} \Leftrightarrow \gamma \leq c(C' d \log(d))^{-1}.\]
    This condition can again be ensured by decreasing $\gamma_0(\delta, \epsilon)$ as needed. Finally, let 
    \[ \beta = \frac{\tilde\lambda}{\bar\lambda - Ck(\epsilon+q)}, \]
    where $\tilde \lambda$ is the one of Proposition~\ref{2405:prop:app:low_dim_exp}. Then $\beta < 1$ and we have at $t = \tau_\eta$
    \[ \norm{\vec{u}_{\tau_\eta} - \bar\Pi^\star \vec{u}_{\tau_\eta}} \leq \frac{K}{\sqrt{d}} \exp\left(\beta \log\left(\frac{\eta\sqrt{d}}{\kappa c'}\right)\right) \leq C''(\delta, \epsilon) d^{\frac{1 - \beta}{2}}. \]
\end{proof}

\section{Proof of the Main Results}
\label{2405:sec:app:proofs}

\subsection{Proof of Theorem \ref{2405:thm:single_index_recovery}}

We formalize here the proof sketch of Theorem \ref{2405:thm:single_index_recovery}.  We track the dynamics of $\vec{w}_t$ through the following sufficient statistic:
\begin{equation}
  m_t := \langle \vec{w}_t, \vec{w}^\star \rangle.
\end{equation}
The dynamics of $m_t$ fall under the framework of Section~\ref{2405:sec:app:framework}, by letting
\[ \Phi(\vec{w}) = -\nabla_{\vec w}^\bot\cL\left(f(\vec z; \tilde{\vec{w}}(\rho), a_0),  y \right). \]
Indeed, for the spherical gradient,
\begin{align*}
  \nabla_{\vec{w}}^\bot \cL(f(\vec{z}; \vec{w}, a_0), y) = a_0 y \sigma'(\langle \vec{w}, z \rangle) \cdot \left(\vec{z} - \langle \vec{z}, \vec{w} \rangle \vec{w} \right).
\end{align*},
hence plugging in the expression for $\tilde{\vec{w}}(\rho)$
\begin{equation}\label{2405:eq:app:phi_expression}
    \Phi(\vec{w}) = h^\star(\langle \vec{w}^\star, \vec{z} \rangle) \sigma'\left(\langle \vec{w}, \vec{z} \rangle + \rho h^\star(\langle \vec{w}^\star, \vec{z} \rangle) \sigma'(\langle \vec{w}, \vec{z} \rangle) \cdot \norm{\vec{z}}^2\right) \left(\vec{z} - \langle \vec{z}, \vec{w} \rangle \vec{w} \right).
\end{equation} 
Assumptions  \ref{2405:assump:poly_growth} and \ref{2405:assump:init} on $\sigma$ and $h^\star$ then directly imply Assumptions \ref{2405:assump:app:initialization} and \ref{2405:assump:app:concentration} in the Appendix. As a result, we only need to study the quantity $\Psi(\vec{w}) = \Ea{\Phi(\vec{w})}$.
We show the following:
\begin{lemma}\label{2405:lem:app:psi_low_rank}
    The function $\Psi$ is approximately one-dimensional: there exists a function $\phi$ such that for any $\vec{w} \in \mathbb{S}^{d-1}$, we have
    \[ |\langle \vec{w}^\star, \Psi(\vec{w}) \rangle- a_0 \phi(\vec{w}, \vec{w}^\star)| \leq \frac{C}d \]
\end{lemma}

\begin{proof}
Define the following quantities:
\begin{align}
  &\Xi(x, x^\star, x^\bot) = h^\star(x^\star) \sigma'(x + a_0\rho_0 \sigma'(x) h^\star(x^\star)) x^\bot \\
  &\phi(m) = \E[\Xi(x, x^\star, x^\bot)] \  \text{for} \  \dbinom{x}{x^\star} \sim  \cN\left( \vec{0},
    \begin{pmatrix}
      1 & m \\
      m & 1
  \end{pmatrix} \right) \  \text{and} \  x^\bot = x^\star - m x.\label{2405:eq:app:def_phi}
\end{align}

  For simplicity, we let $u = \langle \vec{w}, \vec{z} \rangle + a_0\rho_0 h^\star(\langle \vec{w}^\star, \vec{z} \rangle)$. Doing a Taylor expansion of $\sigma'$ in the second order, we have
  \[
    \resizebox{\linewidth}{!}{$\displaystyle
      \begin{aligned}
        \sigma'\left(\langle \vec{w}, \vec{z} \rangle + \rho h^\star(\langle \vec{w}^\star, \vec{z} \rangle) \sigma'(\langle \vec{w}, \vec{z} \rangle) \cdot \norm{\vec{z}}^2\right) = &\sigma'(u)
        + \sigma''(u) \rho h^\star(\langle \vec{w}^\star, \vec{z} \rangle) \sigma'(\langle \vec{w}, \vec{z} \rangle) \cdot  \left(\norm{\vec{z}}^2 - d \right) \\
        &+ R(u) \rho^2 h^\star(\langle \vec{w}^\star, \vec{z} \rangle)^2 \sigma'(\langle \vec{w}, \vec{z} \rangle)^2 \cdot  \left(\norm{\vec{z}}^2 - d \right)^2 ,
      \end{aligned}
    $}
  \]
  where $R(u)$ is the Taylor remainder of $\sigma'$ at second order around $u$.

  The distribution of the triplet $(\langle \vec{w}, \vec{z} \rangle, \langle \vec{w^\star}, \vec{z} \rangle, \langle \vec{w^\star}, \vec{z}^\bot \rangle)$ is the same as the one of $(x, x^\star, x^\bot)$, hence
  \[
    \resizebox{\linewidth}{!}{$\displaystyle
      \begin{aligned}
        \Ea{\langle \Phi(\vec{z}), \vec{w}^\star \rangle} - a_0\phi(m) = &\underbrace{\rho\Ea{\sigma''(u)  h^\star(\langle \vec{w}^\star, \vec{z} \rangle)^2 \sigma'(\langle \vec{w}, \vec{z} \rangle) \cdot  \left(\norm{\vec{z}}^2 - d \right) \langle \vec{w}^\star, \vec{z}^\bot \rangle}}_{(1)} \\
        &+ \underbrace{\rho^2 \Ea{R(u) h^\star(\langle \vec{w}^\star, \vec{z} \rangle)^3 \sigma'(\langle \vec{w}, \vec{z} \rangle)^2 \cdot  \left(\norm{\vec{z}}^2 - d \right)^2  \langle \vec{w}^\star, \vec{z}^\bot \rangle}}_{(2)}.
      \end{aligned}
    $}
  \]
  We first bound (2). From the Hölder inequality (since $\frac12 + \frac18 + \frac18 + \frac18 + \frac18 = 1$),
  \begin{align*}
    (2) &\leq \rho^2 \Ea{R(u)^8}^{\frac18}\Ea{h^\star(x^\star)^8}^{\frac18}\Ea{\sigma'(x)^8}^{\frac18}\Ea{\langle \vec{w}^\star, \vec{z}^\bot \rangle^8}^{\frac18}\Ea{\left(\norm{\vec{z}}^2 - d \right)^4}^{\frac12}
  \end{align*}
  Since $\sigma$ is three times differentiable almost everywhere, and from the polynomial bounds of Assumption~\ref{2405:assump:derivatives}, the first four expectations are bounded by an absolute constant. The last term is the centered fourth moment of a $\chi^2(d)$ distribution, which is equal to $12d(d+4)$. As a result, using that $\rho^2 = O(d^{-2})$,
  \[ (2) \leq Cd^{-1}. \]
  For the first term, we let $\vec{z} = \vec{z}_\parallel + \vec{z}'$, where $\vec{z}_\parallel$ is the projection of $\vec{z}$ on the span of $\vec{w}$ and $\vec{w}^\star$. We can then write
  \begin{align*} 
    (1)& = \underbrace{\Ea{\sigma''(u) \rho  h^\star(\langle \vec{w}^\star, \vec{z} \rangle)^2 \sigma'(\langle \vec{w}, \vec{z} \rangle) \cdot  \norm{\vec{z}_\parallel}^2 \cdot \langle \vec{w}^\star, \vec{z}^\bot \rangle}}_{(1')} \\
    &\quad+ \underbrace{\Ea{\sigma''(u) \rho  h^\star(\langle \vec{w}^\star, \vec{z} \rangle)^2 \sigma'(\langle \vec{w}, \vec{z} \rangle) \cdot  \left(\norm{\vec{z'}}^2 - d \right) \langle \vec{w}^\star, \vec{z}^\bot \rangle}}_{(1'')}.
  \end{align*}
  Up to the factor of $\rho$, the term $(1')$ is the expectation of a polynomially bounded function of a $2$-dimensional Gaussian variable, hence
  \[ (1') \leq C \rho \leq C' d^{-1}. \]
  In $(1'')$, the vector $\vec{z'}$ is independent from $\langle \vec{w}, \vec{z} \rangle$ and $\langle \vec{w}^\star, \vec{z} \rangle$, so the expectation factorizes as
  \begin{align*}
    |(1'')| &= \left|\Ea{\sigma''(u) \rho  h^\star(\langle \vec{w}^\star, \vec{z} \rangle)^2 \sigma'(\langle \vec{w}, \vec{z} \rangle) \langle \vec{w}^\star, \vec{z}^\bot \rangle }\cdot  \Ea{\left(\norm{\vec{z'}}^2 - d \right) }\right| \\
    &\leq C \rho |d - 2 - d| \leq C' d^{-1},
  \end{align*}
  where we used that $\vec{z}'$ is a $(d-2)$-dimensional Gaussian variable. This concludes the proof.
\end{proof}

Lemma \ref{2405:lem:app:psi_low_rank} implies that Assumption~\ref{2405:assump:app:low_dim} holds for $\Psi$ and any choice of $q$, as long as $d$ is large enough.  As a result, the proof of Theorem~\ref{2405:thm:single_index_recovery} hinges on the following facts about $\phi$:
\begin{proposition}\label{2405:prop:app:phi_phip_nnz}
  For almost any choice of $\rho_0$ (w.r.t to the Lebesgue measure),
  \begin{itemize}
    \item if $\ell_p^\star = 1$, then $\phi(0) \neq 0$;
    \item if $\ell_p^\star = 2$, then $\phi'(0) \neq 0$.
  \end{itemize}
\end{proposition}

\begin{proof}[of Theorem \ref{2405:thm:single_index_recovery}]
    Our goal is to apply Proposition~\ref{2405:prop:app:hitting_times} to the function $\phi$. If $\ell_p^\star = 1$, then $\phi(0) \neq 0$ by the previous lemma, and with a uniform initialization of $\vec{w}^\star$ we have $m_0 \phi(0) > 0$ with probability $1/2$.  Combining this with the $(1- \delta)$ probability event of Proposition \ref{2405:prop:app:hitting_times} for small enough $\epsilon$, with probability $1 - \delta$, we have $t_\eta \leq C(\delta)d^{-1}$ with probability $1/2 - \delta$.
    We now move to the proof of the case where $\ell_p^\star =1$. In this case, Proposition~\ref{2405:prop:app:hitting_times} implies the desired result as soon as $\phi'(0) > 0$ (which corresponds in the $k = 1$ setting to the only eigenvalue of $d\phi(0)$ being positive). Since the Lebesgue measure is invariant w.r.t multiplication by $\pm 1$, the factor $a_0 \rho_0$ in the definition of $\phi'(0)$ has the same distribution as $\rho_0$, hence we can write $\phi'(0) = a'_0 |\phi'(0)|$ where $a'_0$ is also uniform in $\{-1, 1\}$. The event $\phi'(0) > 0$ therefore also has probability $1/2$, and we conclude as before.
\end{proof}

\subsubsection*{Analysis of $\phi$: proof of Proposition~\ref{2405:prop:app:phi_phip_nnz}}

The main ingredient here is to write $\phi(m) = \phi(m; \rho_0)$ and differentiate w.r.t $\rho_0$. For simplicity, we take $a_0 = 1$ in the definition of $\phi$. We show the following:
\begin{lemma}\label{2405:lem:app:phi_diff}
  Under Assumption \ref{2405:assump:derivatives}, the function $\phi(m; \rho_0)$ is analytic in $\rho_0$. Further, we have
  \begin{equation}
    \left.\frac{\partial^k \phi(m; \rho_0)}{\partial \rho_0^k}\right|_{\rho_0 = 0} = \Ea{h^\star(x^\star)^{k+1} x^\bot \sigma^{(k+1)}(x) \sigma'(x)^k}
  \end{equation}
  where $x, x^\star$ follow the same distribution as in eq.~\eqref{2405:eq:app:def_phi}.
\end{lemma}

\begin{proof}
  The first statement stems from the analyticity of $\sigma$, and the fact that it is a property conserved through integration. For the second, we differentiate inside the expectation, to find
  \[ \frac{\partial \phi}{\partial \rho_0} = \Ea{\sigma'(x)h^\star(x^\star)^2 \sigma''(x + \rho_0 \sigma'(x)h^\star(x^\star)) x^\bot}; \]
  the result follows by induction and taking $\rho_0 = 0$ at the end.
\end{proof}

We define the following quantities related to Assumption~\ref{2405:assump:init}:
\begin{align*}
  u_{0}^{(k)} &= \Ea{\sigma^{(k)}(x) \sigma'(x)^{k-1}}  &u_{1}^{(k)} &= \Ea{x \sigma^{(k)}(x) \sigma'(x)^{k-1}} \\
  v_{0}^{(k)} &= \Ea{x^\star h^\star(x^\star)^k} & v_{1}^{(k)} &= \Ea{[(x^\star)^2-1] h^\star(x^\star)^k} \\
  \Delta^{(k)} &= \Ea{h^\star(x^\star)^k}
\end{align*}
Then, by Proposition 11.37 from \cite{odonnell2014analysis}, and using $x^\bot = x^\star - mx$:
\[
  \resizebox{\linewidth}{!}{$\displaystyle
    \begin{aligned}
      \left.\frac{\partial^k \phi(m; \rho_0)}{\partial \rho_0^k}\right|_{\rho_0 = 0} &= \Ea{h^\star(x^\star)^{k+1} x^\star \sigma^{(k+1)}(x) \sigma'(x)^k} - m \cdot \Ea{h^\star(x^\star)^{k+1} x \sigma^{(k+1)}(x) \sigma'(x)^k} \\
      &= \left(u_{0}^{(k)}v_{0}^{(k)} + u_{1}^{(k)}(v_{1}^{(k)} + \Delta^{(k)})\cdot m + o(m)\right) - m \left(u_1^{(k)} \Delta^{(k)} + o(1)\right) \\
      &= u_{0}^{(k)}v_{0}^{(k)} + u_{1}^{(k)}v_{1}^{(k)}\cdot m + o(m)
    \end{aligned}
  $}
\]
This allows us to show Proposition \ref{2405:prop:app:phi_phip_nnz}:
\begin{proof}
  Assume first that $\ell_p^\star = 1$. Then $\phi(0; \rho_0)$ is an analytic function of $\rho_0$ with
  \[ \left.\frac{\partial^k \phi(0; \rho_0)}{\partial \rho_0^k}\right|_{\rho_0 = 0} = u_{0}^{(k)}v_{0}^{(k)}. \]
  Since $\ell_p^\star = 1$, there exists a $k \in \mathbb{N}$ such that $v_{0}^{(k)} \neq 0$, and by Assumption~\ref{2405:assump:derivatives} the coefficient $u_{0}^{(k)}$ is nonzero for every $k$. As a result, $\phi(0;\rho_0)$ is an analytic and non-identically zero function of $\rho_0$, so it is non-zero for almost every choice of $\rho_0$, as requested. The case $k_p = 2$ is done in a similar way, noting that this time
  \[ \left.\frac{\partial^k \phi'(0; \rho_0)}{\partial \rho_0^k}\right|_{\rho_0 = 0} = u_{1}^{(k)}v_{1}^{(k)}. \]
\end{proof}

\subsection{Proof of Theorem \ref{2405:thm:multi_index_recovery}}

The proof of Theorem~\ref{2405:thm:multi_index_recovery} follows in the same way as the one of Theorem~\ref{2405:thm:single_index_recovery}. Recall the definition of $\Phi$ in eq.~\eqref{2405:eq:app:phi_expression}, and define
\begin{align}
    &\Xi(z, \vec{z}^\star, \vec{z}^\bot) = h^\star(\vec{z}^\star) \sigma'(z + a_0\rho_0\sigma'(z)h^\star(\vec{z}^\star)) \vec{z}^\bot \\
    &\psi(\vec{p}) = \Ea{\Xi(z, \vec{z}^\star, \vec{z}^\bot)} \  \text{for} \  \dbinom{z}{\vec{z}^\star} \sim  \cN\left( \vec{0},
    \begin{pmatrix}
      1 & \vec{p} \\
      \vec{p} & I_k
  \end{pmatrix} \right) \  \text{and} \  \vec{z}^\bot = \vec{z}^\star - z\vec{p}.\label{2405:eq:app:def_psi}
\end{align}
Again, the variables $(z, \vec{z}^\star)$ correspond to $\langle \vec{w}, \vec{z} \rangle, W^\star\vec{z}$, and the vector $\vec{p}$ to $W^\star\vec{w}$. As before, the only difference between $\Ea{\Phi(\vec{w})}$ and $\psi(W^\star \vec{w})$ lies in replacing $\norm{\vec{z}}$ with $d$, and the proof of Lemma~\ref{2405:lem:app:psi_low_rank} implies the following: for any $\vec{w}\in \mathbb{S}^{d-1}$,
\[ \norm{W^\star\Psi(\vec{w}) - \psi(W^\star \vec{w})} \leq \frac Cd. \]
As a result, the function $\Psi$ is approximately $k$-dimensional for any choice of $q$ as $d \to \infty$. 

We can view the vector $\psi(\vec{p})$ as a random variable depending on the choice of $a_0, \rho_0$. Since $a_0$ is a sign, we will only consider the dependency of $\psi$ in $\rho_0$. As before, let $u_i^{(k)}$ be the $i$-th Hermite coefficient of $z \mapsto \sigma^{(k+1)}(z) \sigma'(z)^k$, and $C_i^{(k)}$ the $i$-th Hermite tensor of $h^\star$ \citep{grad1949note} -- in particular, $C_i^{(k)}$ is a tensor of order $i$.
We have the analogous result to Lemma~\ref{2405:lem:app:phi_diff}:

\begin{lemma}
    Under Assumption~\ref{2405:assump:derivatives}, the function $\psi(\vec{p}; \rho_0)$ is analytic in $\rho_0$, with
    \[  \left.\frac{\partial^k \psi(\vec{p}; \rho_0)}{\partial \rho_0^k}\right|_{\rho_0 = 0} = u_0^{(k)}C_1^{(k)} + u_1^{(k)}C_2^{(k)} \vec{p} + o(\norm{\vec{p}}). \]
\end{lemma}

\begin{proof}
    As in Lemma~\ref{2405:lem:app:phi_diff}, we can write $\psi(\vec{0}) = \psi(\vec{0}, \rho_0)$, with
    \[  \left.\frac{\partial^k \psi(\vec{p}; \rho_0)}{\partial \rho_0^k}\right|_{\rho_0 = 0} = \Ea{h^\star(\vec{z}^\star)^{k+1} \vec{z}^\bot \sigma^{(k+1)}(z) \sigma'(z)^k}. \]
     From \cite[Lemma 19]{dandi2023how}, we have
    \[ \Ea{h^\star(\vec{z}^\star)^{k+1} \vec{z}^\star \sigma^{(k+1)}(z) \sigma'(z)^k} = \sum_{i=0}^\infty u_{i+1}^{(k)} \langle \vec{p}^{\otimes i}, C_i^{(k)} \rangle \vec{p} + \sum_{i = 0}^\infty u_{i}^{(k)} C_{i+1}^{(k)} \times_{1, \dots, i} \vec{p}^{\otimes i}. \]
    Since $\vec{z}^\bot = \vec{z} - z\vec{p}$, we find
    \begin{equation}
        \left.\frac{\partial^k \psi(\vec{p}; \rho_0)}{\partial \rho_0^k}\right|_{\rho_0 = 0} = \sum_{i=0}^{\infty} u_i^{(k)} \left( C_{i+1}^{(k)} \times_{1, \dots, i} \vec{p}^{\otimes i} - \langle C_{i+1}^{(k)}, \vec{p}^{\otimes (i+1)} \rangle \vec{p} \right).
    \end{equation}
    For $i \geq 0$, we have $\norm{C_{i+1}^{(k)} \times_{1, \dots, i} \vec{p}^{\otimes i}} = O\left(\norm{\vec{p}}^i\right)$ and $\langle C_{i}^{(k)}, \vec{p}^{\otimes (i)} \rangle = O\left(\norm{\vec{p}}^i\right)$. The only terms in the sum which are not of order $\norm{\vec{p}}^2$ are the first two terms in the left sum, which yields the statement of the Lemma.
\end{proof}

We are now ready to show the first part of Theorem~\ref{2405:thm:multi_index_recovery}. We begin with a characterization of $U^\star$ whenever $\ell_p^\star = 1$:
\begin{lemma}
    If $\ell_p^\star = 1$, then
    \[ U^\star = \operatorname{span}\left(\left\{C_1^{(k)}\right\}_{k \geq 1}\right). \]
    If $\ell_p^\star = 2$, then
    \[ U^\star = \bigoplus_{k\geq 1}\mathrm{Im}\left(C_2^{(k)}\right). \]
\end{lemma}

\begin{proof}
    By the properties of the Hermite tensors, we have for any $v \in V^\star$ and any polynomial $p$
    \[ \Ea{y^k H_1(\langle \vec{v}, \vec{z} \rangle)} = \langle \vec{v}, C_1^{(k)} \rangle \]
    As a result, if $\vec{v} = \sum_{k \geq 0} v_k C_1^{(k)}$, then by letting $p(y) = \sum_{k \geq 0} a_k y^k$ we have
    \[ \Ea{p(y) H_1(\langle \vec{v}, \vec{z} \rangle)} = \norm{\vec{v}}^2 \neq 0.\]
    Conversely, if $\vec{v}$ is orthogonal to every vector $C_1^{(k)}$, then $\Ea{y^k H_1(\langle \vec{v}, \vec{z} \rangle)} = 0$ for all $k \geq 0$, which implies that $\vec{v} \notin U^\star$.

    The case $\ell_p^\star = 2$ is similar, since this time
    \[ \Ea{y^k H_2(\langle \vec{v}, \vec{z} \rangle)} = \vec{v}^\top C_2^{(k)} \vec{v}. \]
    As a result, for fixed $k$, the largest vector space $U_k$ such that $\Ea{y^k H_2(\langle \vec{v}, \vec{z} \rangle)}$ for all $\vec{v}\in U_k$ is $\ker(C_2^{(k)}) = \mathrm{Im}(C_2^{(k)})^\bot$. The result ensues from $U^\star = \bigcap_{k \geq 1} U_k^\bot$.
\end{proof}

Using Assumption~\ref{2405:assump:derivatives}, the following corollary ensues:
\begin{corollary}
    If $\ell_p^\star = 1$, then for any $\vec{v}\in U^\star$ and almost every $\rho_0$,
    \[ \langle \vec{v}, \psi(\vec{0}) \rangle \neq 0. \]
    As a result, with probability one, the set $\{ \psi(\vec{0}; \rho_{0, i}) \}_{i \in [p]}$ spans $U^\star$ as soon as $p \geq k$.
    
    If $\ell_p^\star = 2$, then with probability at least $1/2$ on the choice of $\rho_0, a_0$, $d\psi(\vec{0})$ is symmetric and has a strictly positive eigenvalue.
\end{corollary}

\begin{proof}
    Let $\vec{v}\in U^\star$, and $k \geq 1$ such that $\langle \vec{v}, C_1^{(k)} \rangle \neq 0$. Then 
    \[\left.\frac{\partial^k \langle \vec{v}, \psi(\vec{0}; \rho_0) \rangle}{\partial \rho_0^k}\right|_{\rho_0 = 0} = u_0^{(k)} \langle \vec{v}, C_1^{(k)} \rangle \neq 0, \]
    so the function $\rho_0 \mapsto \langle \vec{v}, \psi(\vec{0}; \rho_0) \rangle$ is not the zero function. Since it is also analytic, it is nonzero for almost every choice of $\rho$.
    
    We now prove that if $p \leq \dim(U^\star)$, then the family $\{ \psi(\vec{0}; \rho_{0, i}) \}_{i \in [p]}$ is free, which implies our second result. Assume by recursion that $\{ \psi(\vec{0}; \rho_{0, i}) \}_{i \in [p-1]}$ is free, and let $\vec{v}_p \in U^\star$ orthogonal to this family. Then with probability one, $\langle \vec{v}_p, \psi(\vec{0}; \rho_{0, p}) \rangle \neq 0$, so $\psi(\vec{0}; \rho_{0, p})$ is independent from $\{ \psi(\vec{0}; \rho_{0, i}) \}_{i \in [p-1]}$.

    Finally, if $\ell_p^\star = 2$, then by symmetry of the distribution of $\rho_0$ we can write $\psi(\vec{p}; a_0, \rho_0) = a_0 \psi(\vec{p}; 1, \rho')$ with $\rho' = a_0 \rho_0$ chosen independently from $a_0$. Writing
    \[ d\psi(\vec{0}; \rho') = \sum_{k \geq 1} (\rho')^k u_1^{(k)} C_2^{(k)}, \]
    we can see that $d\psi(\vec{0})$ is a symmetric matrix, which is non-zero with probability one on $\rho'$. Since either $d\psi(\vec{0})$ or $-d\psi(\vec{0})$ has a positive eigenvalue, $a_0d\psi(\vec{0})$ does with probability $1/2$.
\end{proof}

We are now ready to prove Theorem~\ref{2405:thm:multi_index_recovery}. 
\begin{proof}[of Theorem~\ref{2405:thm:multi_index_recovery}]
    We begin with the case $\ell_p^\star = 1$. Fix a $p = p(\delta)$ to be chosen later, and let $\psi_i(\vec{p}) = \psi(\vec{p}; a_{0, i}, \rho_{0, i})$ for $i \in [p]$. With probability $1/2$ on $\vec{w}_{0, i}$, we have $\langle \vec{w}_{0, i}, \psi_i(\vec{0}) \rangle$; let $p_+$ be the number of such neurons. For $p = O(\log(1/\delta))$, we have $p_+ > \max(p/4, k)$ with probability at least $1 - \frac\delta2$.

    In particular, with probability one, the set $\{ \psi_i(\vec{0}) \}_{i\in [p_+]}$ spans $U^\star$. Let $\epsilon > 0$ to be chosen later, $\eta_i(\epsilon)$ adapted to $\psi_i$ in Proposition~\ref{2405:prop:app:hitting_times}, $\tau_i$ the hitting time for $\eta_i$ in the $i$-th neuron. Consider the time $\tau = \min_i \tau_i$; applying Proposition \ref{2405:prop:app:hitting_times} to $\delta' = \delta/p$, we have
    \[ \tau \leq \max_{i \in [p_+]} \frac{\eta_i(1 + \sqrt{k}\epsilon)}{\gamma_0 \norm{\psi_i(\vec{0})}} d \leq C(\delta) d. \]
    However, we also have $\tau > c(\delta) \gamma^{-1}$, and hence from Proposition~\ref{2405:prop:app:low_dim_linear}, for any $i \in [p_+]$
    \[ \norm{\vec{w}_{t, i}} \geq \gamma \tau \norm{\psi_i(\vec{0})} - \sqrt{k}\eta_i\epsilon \geq c'(\delta) \]
    for small enough $\epsilon$. This proves the first part of the theorem. Now, consider the vectors
    \[ \vec{u}_i = \frac{\Pi^\star \vec{w}_{\tau, i}}{\norm{\Pi^\star \vec{w}_{\tau, i}}}; \]
    by the same argument as in Proposition~\ref{2405:prop:app:hitting_times}, for any $i \in [p_+]$,
    \[ \norm{\vec{u}_i - \frac{\psi_i(\vec{0})}{\norm{\psi_i(\vec{0})}}} \leq \sqrt{k}\epsilon. \]
    Since the $\{\psi_i(\vec{0})\}_{i\in[p_+]}$ span $U^\star$ with probability one, so do the $\{\vec{u}_i\}_{i \in [p_+]}$ for a small enough choice of $\epsilon$.

    For the case $\ell_p^\star = 2$, we define the $p_+$ ``good'' neurons as those with $d\psi_i(\vec{0})$ having a positive eigenvalue. Since this happens with probability at least $1/2$ independently for each neuron, the bounds on $p(\delta)$ from earlier still apply.
    For those neurons, and small enough $\epsilon$, we define the $\eta_i(\epsilon)$ as those in Proposition~\ref{2405:prop:app:hitting_times}, and $\tau = \max_i \tau_i$. Then by Proposition~\ref{2405:prop:app:hitting_times}, we have $\tau \leq C(\delta)d\log(d)^2$ for some constant $C(\delta)$.

    Now, fix a vector $i \in [p_+]$. We define the $j$-th ``excursion'' to be the $j$-th sequence of times $\tau_i^{(j)}, \dots, \tau_i^{(j)} + e_i^{(j)}$ such that $\norm{\Pi^\star\vec{w}_{t, i}} \geq \eta_i$. Then:
    \begin{itemize}
        \item  With probability at least $1 - e^{-c\gamma}$, we have $\norm{\vec{w}_{t+1, i} - \vec{w}_{t, i}} \leq C \sqrt{\gamma}$, and thus for any excursion $j$ we have $\norm{\Pi^\star\vec{w}_{\tau_i^{(j)} + e_i^{(j)} +1, i}} \geq \eta_i/2$.
        \item between two excursions, the results of Proposition~\ref{2405:prop:app:low_dim_exp} apply, and thus the norm of $\norm{\Pi^\star \vec{w}_{t, i}}$ increases exponentially.
    \end{itemize}
   As a result, for $\tau_i \leq t \leq \tau$, we have $\norm{\Pi^\star \vec{w}_{t, i}} > \eta_i / 2$, which finishes the proof of Theorem~\ref{2405:thm:multi_index_recovery}.
\end{proof}

\section{Additional proofs}

\subsection{Proof of Lemma~\ref{2405:lem:bias_activation}}

Let $\sigma$ be an analytic function which is not a polynomial. Then for any $k \geq 0$, the function $\sigma^{(k)}$ is also analytic and non-identically zero, and hence so is the function $f_n = \sigma^{(k)} \cdot (\sigma')^k$. Lemma \ref{2405:lem:bias_activation} thus ensues from the following result:
\begin{lemma}
  Let $f$ be an analytic and non-identically zero function. Then
  \[ \Ea{f(x+b)} \neq 0 \quad \text{and} \quad \Ea{x f(x+b)} \neq 0 \]
  for almost every $b$ under the Lebesgue measure.
\end{lemma}
\begin{proof}
  The function $\psi(b) = \Ea{f(x+b)}$ is analytic in $b$, and we have
  \[ \psi^{(k)}(0) = \Ea{f^{(k)}(x)} = c_k u_k(f), \]
  where $u_k(f)$ is the $k$-th Hermite coefficient of $f$ and $c_k$ is a nonzero absolute constant. Since $f$ is nonzero, at least one of its Hermite coefficients is nonzero, and $\psi$ is a non-identically zero analytic function. As a result, $\psi(b) \neq 0$ for almost every choice of $b$.

  The other case is handled as the first one, noting that by Stein's lemma
  \[ \Ea{x f(x+b)} = \Ea{f'(x+b)}. \]
\end{proof}
The above shows that for fixed $n \in \mathbb{N}$, the set
\[
  \resizebox{\linewidth}{!}{$\displaystyle
    \mathcal{B}_n = \left\{ b \in \R \ : \ \Ea{\sigma^{(n)}(x+b) \sigma'(x+b)^{n-1}} = 0 \quad \text{or} \quad \Ea{x \sigma^{(n)}(x+b) \sigma'(x+b)^{n-1}} = 0  \right\}
  $}
\]
has measure $0$. Since there is a countable number of such subsets, the set
\[ \mathcal{B} = \bigcup_{n \in \mathbb{N}} \mathcal{B}_n \]
also has measure zero, which is equivalent to the statement of Lemma~\ref{2405:lem:bias_activation}.

\subsection{Proof of Theorem~\ref{2405:thm:poly_pge}}

We decompose the polynomial $h^\star$ as an even and odd part $e(x)$ and $o(x)$, with degrees $d_o$ and $d_e$. We first assume that $o$ is non-zero, and that the leading coefficient of $o$ is positive. Then, for odd $m \geq 0$, we have
\[ \Ea{xh^\star(x)^m} = \sum_{k = 0}^m \dbinom{m}{k} \Ea{x e(x)^k o(x)^{m-k}} \]
Each term in the above sum is zero if $k$ is odd, we focus on the terms where $k$ is even. There exist constants $A, \varepsilon > 0$ such that if $|x| \geq A$,
\[ |e(x)| \geq \varepsilon |x|^{d_e} \quad \text{and} \quad |o(x)| \geq \varepsilon |x|^{d_o}, \]
and we let
\[ B = \sup_{x \in  [-A, A]} |e(x)| \vee \sup_{x \in  [-A, A]} |o(x)| \]
As a result, when $k$ is even, both $e(x)^k$ and $xo(x)^{m-k}$ are even polynomials with positive leading coefficients, so
\begin{align*}
  \Ea{x e(x)^k o(x)^{m-k}} &= \Ea{x e(x)^k o(x)^{m-k} \mathbf{1}_{x\in [-A, A]}} + \Ea{x e(x)^k o(x)^{m-k} \mathbf{1}_{x\notin [-A, A]}} \\
  &\geq \varepsilon^2\Ea{ x^{kd_e + (m-k)d_o+1} \mathbf{1}_{x\notin [-A, A]}} - AB^m
\end{align*}
Let $d(m, k) = kd_e + (m-k)d_o+1$. Then,
\begin{align*}
  \Ea{x e(x)^k o(x)^{m-k}} \geq \varepsilon^2\Ea{x^{d(m, k)}} -  A^{d(m, k)} - AB^m.
\end{align*}
Going back to the sum, and with the crude bound $\dbinom{m}{k} \leq 2^m$,
\begin{align*}
  \Ea{xh^\star(x)^m} &\geq \varepsilon^2\Ea{x^{d(m, 0)}} - \sum_{k=1}^m 2^m\left(A^{d(m, k)} + AB^m\right) \\
  &\geq \Ea{x^{m+1}} - C^m
\end{align*}
Since the Gaussian moments grow faster than any power of $m$, this last expression is non-zero for a large enough choice of $m$.

Finally, if $o$ is the zero polynomial, we can do the same reasoning with $e(x)^m$ to get
\[ \Ea{(x^2-1) e(x)^m} = \varepsilon^2(\Ea{x^{m d_e+2}} - \Ea{x^{m d_e}}) - C^m = \varepsilon^2 md_e \Ea{x^{m d_e}} - C^m,\]
and we conclude as before.

\subsection{Hardness of sign functions}
We show in the appendix the following statement:
\begin{proposition}\label{2405:prop:app:sign_hard}
  Let $m \in \mathbb{N}$, and
  \[ h^\star(\vec{x}) = \mathrm{sign}(x_1\dots x_m) \]
  Then the function $h^\star$ has Information and Generative Exponents $\ell = \ell^\star = m$.
\end{proposition}

We first show the lower bound. For any vector $\vec{v}\in \R^m$, and $k < m$, $H_k(\langle \vec{v}, \vec{x} \rangle)$ is a polynomial in $z$ of degree $k$, and therefore each monomial term is missing at least one variable. We show that this implies the lower bound of Proposition~\ref{2405:prop:app:sign_hard} through the following lemma:
\begin{lemma}
  Let $f: \{-1, 1\} \to \R$ be an arbitrary function, and $g: \R^m \to \R$ a function which is independent from $x_m$. Then
  \[ \Ea{f(h^\star(\vec{x})) g(\vec{x})} = \frac{f(1)+f(-1)}{2}\Ea{g(\vec{x})}\]
\end{lemma}
\begin{proof}
  We simply write
  \begin{align*}
    \Ea{f(h^\star(\vec{x})) g(\vec{x})}
    &= \Ea{\Ea{f(h^\star(\vec{x})) g(\vec{x}) \mid x_1, \dots, x_{m-1}}}\\
    &= \Ea{g(\vec{x}) \Ea{f(h^\star(\vec{x})) \mid x_1, \dots, x_{m-1}}} \\
    &= \Ea{g(\vec{x}) \frac{f(1)+f(-1)}{2}} \\
    &=  \frac{f(1)+f(-1)}{2}\Ea{g(\vec{x})},
  \end{align*}
  where at the second line we used that $g$ only depends on $x_1, \dots, x_{m-1}$ and at the third line the addition of $x_m$ randomly flips the sign of $h^\star$.
\end{proof}

Using this property on every monomial in $H_k(\langle \vec{v}, \vec{x} \rangle)$ and summing back, we have for any function $f$
\begin{equation}
  \Ea{f(h^\star(x)) H_k(\langle \vec{v}, \vec{x} \rangle)} = \frac{f(1)+f(-1)}{2}\Ea{H_k(\langle \vec{v}, \vec{x} \rangle)}
\end{equation}
But the latter is zero when $k\geq 1$ by orthogonality of the Hermite polynomials. This shows that $\ell_{\vec{v}}^\star \geq m$ for all $\vec{v}\in\R^k$, and therefore $m \leq \ell^\star \leq \ell$.

For the upper bound, let $\vec{v} = \vec{e}_1 + \dots + \vec{e_m}$. Then
\[ H_k(\langle \vec{v}, \vec{x} \rangle) = m!   x_1\dots x_m + Q(\vec{x}), \]
where $Q$ is a polynomial where each monomial has a missing variable. As a result,
\begin{align*}
  \Ea{h^\star(\vec{x}) H_k(\langle \vec{v}, \vec{x} \rangle)} &= m!  \Ea{h^\star(\vec{x}) x_1\dots x_m} \\
  &= m!  \Ea{|x_1\dots x_m|} \\
  &= m! \left( \sqrt{\frac2\pi} \right)^m
\end{align*}
which is a non-zero value, hence $\ell^\star \leq \ell \leq m$.
